\subsection{Statement of the sharp profile estimates}
\label{hardgap:statement}

For an admissible edge \(n\to m<n\) in \([0,N]\), recall that
\(2n-m\le N\) and
\[
 c_g(m,n)=g(m)-\min_{[m,n]}g.
\]
All chains in this section are finite, decreasing, and have no prescribed
intermediate endpoint. Define
\begin{align}
 b_{\mathrm c}(a)&=\frac{1+2a}{4+2a},
       \label{hardgap:transition}\\
 R(a,b)&=4(1+a)b^2-4(a+2)b+1+3a,
       \label{hardgap:R}
\end{align}
\begin{subequations}\label{hardgap:costs}
\begin{align}
 C_1(a,b)&=\frac{b-a}{1+2b},
       \label{hardgap:C1}\\
 C_2(a,b)&=\frac{(b-a)(2+a-2b-4ab)}{(1-a)(1+2b)^2}.
       \label{hardgap:C2}
\end{align}
\end{subequations}
The additional low-side guard is
\begin{equation}\label{hardgap:guard}
 b\le b_{\mathrm c}(a),\qquad R(a,b)\le0.
\end{equation}

\begin{theorem}[Sharp estimates from hard gaps]\label{hardgap:profile}
Let \(0<a<b\le1/2\), and let \(g:[0,N]\to\mathbb R\) be
1-Lipschitz with \(ax\le g(x)\le bx\). For every
\(0\le n_0<N\), an admissible chain from \(n_0\) to zero exists
with total cost at most
\[
 \begin{cases}
 NC_2(a,b),&b\le b_{\mathrm c}(a)\text{ and }R(a,b)\le0,\\
 NC_1(a,b),&b\ge b_{\mathrm c}(a).
 \end{cases}
\]
Its number of edges is at most
\[
 K\le C\left(1+\log\frac{N}{N-n_0}\right)
\]
with an absolute constant. If a grid of mesh \(T\) contains
\(0,n_0,N\), its endpoints may be chosen on that grid with additional
\(O(KT)\) cost. Uniform additive barrier errors of size \(\epsilon N\)
give additional cost \(O(K\epsilon N+KT)\).

Both bounds are sharp as universal profile estimates in their respective
stated domains: after normalizing \(N=1\), the supremum over admissible
profiles of the limiting optimal cost as \(n_0\uparrow1\) equals the
displayed \(C_i(a,b)\). This is sharpness of the profile model, not a
necessary condition for positive-length distance sets.
\end{theorem}

The previously proved profile estimate handles \(b\ge1/2\). We prove
the theorem by analyzing complementary intervals of the hard-point set,
and by giving explicit nonnegative linear combinations of their constraints.
No numerical optimization is used in the proof.

\subsection{Endpoint normalization}
\label{hardgap:normalization}

Normalize \(N=1\). Replace \(g\) by
\[
 \widetilde g(x)=\min\{g(x),1+a-x\}.
\]
This function is 1-Lipschitz and has the same barriers: the additional
upper bound majorizes \(ax\), since
\(1+a-x-ax=(1+a)(1-x)\ge0\). It has endpoint value
\(\widetilde g(1)=a\). More importantly,
\begin{equation}\label{hardgap:clip}
 c_g(m,n)\le c_{\widetilde g}(m,n)
 \quad(0\le m<n\le1).
\end{equation}
If \(g(m)\le1+a-m\), the lower-endpoint value is unchanged and
the interval minimum can only decrease. Otherwise
\[
 c_{\widetilde g}(m,n)
 \ge(1+a-m)-(1+a-n)=n-m\ge c_g(m,n),
\]
where the last step uses the Lipschitz property. Thus a chain constructed
for the normalized profile gives no greater cost for the original one,
with identical endpoints and length. Henceforth assume \(g(1)=a\).

\subsection{Reversal and the hard-component identity}
\label{hardgap:identity-section}

For \(0<\zeta\le1\), let \(\Phi_\zeta(g)\) be the infimum of
the costs of admissible chains from \(1-\zeta\) to zero. Put
\[
 f(t)=g(1-t)-a,\qquad f(0)=0,\qquad f(1)=-a.
\]
If an original edge has endpoints \(n>m\), write
\(x=1-n\), \(y=1-m\). Its admissibility is exactly
\(y\le2x\). Its original cost becomes
\(f(y)-\min_{[x,y]}f\), whereas the usual total-drop summand is
\(f(x)-\min_{[x,y]}f\). Their difference is \(f(y)-f(x)\).

Let \(T_\zeta(f)\) be the infimum of the usual total drop over
finite partitions of \([\zeta,1]\) whose neighboring endpoints have
ratio at most two. Reversal gives a bijection of these partitions with
the original chains, and telescoping yields
\begin{equation}\label{hardgap:reversal}
 \Phi_\zeta(g)=T_\zeta(f)-g(1-\zeta).
\end{equation}
Let \(T(f)\) denote the infimum of total drop over infinite good
partitions of \([0,1]\), in the sense of Keleti--Shmerkin
\cite[Definition 5.1]{ks}. Extending a finite partition below \(\zeta\)
by successive halvings adds at most \(\zeta\) to its cost. Conversely,
truncate an almost optimal infinite partition at the interval crossing
\(\zeta\), replacing that interval's lower endpoint by \(\zeta\).
The ratio condition is preserved, and the Lipschitz property changes the
crossing cost by at most \(2\zeta\); later costs are nonnegative and
can be discarded. Consequently
\[
 T(f)-\zeta\le T_\zeta(f)\le T(f)+2\zeta.
\]
Equation~\eqref{hardgap:reversal} proves
\begin{equation}\label{hardgap:limiting-cost}
 \Phi_0(g):=\lim_{\zeta\downarrow0}\Phi_\zeta(g)=T(f)-a.
\end{equation}

Suppose for now that \(g\) is piecewise affine with finitely many pieces.
Its hard-point set in the original coordinate is
\[
 H_g=\left\{x\in[0,1]:
           g(x)=\min_{[x,(1+x)/2]}g\right\}.
\]
It is a finite union of closed interval components, possibly degenerate.
Both zero and one belong to it. Under \(t=1-x\), this is exactly the
hard-point set \(\{t:f(t)=\min_{[t/2,t]}f\}\) of \(f\).
On each original component \([p,q]\), the function \(g\) is
nondecreasing. This also follows directly: for sufficiently close
\(x<y\) in a hard interval, hardness at \(x\) gives
\(g(y)\ge g(x)\), and subdivision and continuity give the assertion
throughout the interval.

Keleti--Shmerkin's hard-point formula
\cite[Lemma 5.4]{ks} applies to this piecewise affine function. A component
\([p,q]\) corresponds to \([1-q,1-p]\), whose drop is
\(f(1-q)-f(1-p)=g(q)-g(p)\). Therefore
\begin{equation}\label{hardgap:identity}
 \Phi_0(g)=
 \sum_{[p,q]\text{ component of }H_g}\bigl(g(q)-g(p)\bigr)-a.
\end{equation}
The subtraction of \(a\) is the endpoint correction in
\eqref{hardgap:limiting-cost}; omitting it would change the problem.

\subsection{Gaps, their drops, and two tail estimates}
\label{hardgap:geometry}

\begin{lemma}[A gap ends at its minimum]\label{hardgap:gap}
If \((q,p)\) is a gap between consecutive components of \(H_g\), then
\[
 g(x)\ge g(p)\quad(q<x<p),\qquad g(q)\ge g(p).
\]
Writing \(u=g(q)\), \(v=g(p)\), and \(\delta=u-v\ge0\),
every gap with \(\delta>0\) satisfies
\begin{equation}\label{hardgap:geometric-drop}
 \delta\le p-\frac{1+q}{2}.
\end{equation}
\end{lemma}

\begin{proof}
If \(g(x)<g(p)\) for some \(x\in(q,p)\), choose a minimum point
\(z\) on \([x,p]\). Then \(z<p\), and \(g(z)<g(p)\).
On \([z,p]\) the profile is at least \(g(z)\). Any part of
\([z,(1+z)/2]\) lying after \(p\) is contained in
\([p,(1+p)/2]\), where hardness of \(p\) gives
\(g\ge g(p)>g(z)\). Thus \(z\) is hard, a contradiction.
Continuity gives \(g(q)\ge g(p)\).

If \(\delta>0\), hardness of \(q\) forces
\(p>(1+q)/2\), and \(g((1+q)/2)\ge u\). The Lipschitz
bound between \((1+q)/2\) and \(p\) gives
\eqref{hardgap:geometric-drop}.
\end{proof}

Telescoping all hard-component increases and the intervening gap changes,
using \(g(0)=0\), \(g(1)=a\), turns \eqref{hardgap:identity} into
\begin{equation}\label{hardgap:drop-sum}
 \Phi_0(g)=\sum_{\text{gaps}}\delta.
\end{equation}
All terms are nonnegative. We may omit gaps with zero drop. The retained
positive gaps remain ordered and disjoint, and all Lipschitz and barrier
constraints at their endpoints remain valid.

\begin{lemma}[Two estimates after a specified endpoint]\label{hardgap:tail}
Suppose \(g(p)=v\), and let \(T\) be the sum of the drops of
positive gaps contained in \([p,1]\). Then
\begin{equation}\label{hardgap:tails}
 T\le\frac{1-p}{2},\qquad
 T\le\frac{1-p+v-a}{2}.
\end{equation}
In particular, for \(0\le\chi\le1/2\),
\begin{equation}\label{hardgap:blended-tail}
 T\le\frac{1-p}{2}+\chi(v-a).
\end{equation}
\end{lemma}

\begin{proof}
For a positive gap \((q_j,p_j)\),
\[
 \delta_j\le p_j-\frac{1+q_j}{2}
 =\frac{p_j-q_j}{2}-\frac{1-p_j}{2}
 \le\frac{p_j-q_j}{2}.
\]
Sum over the disjoint intervals to obtain the first bound. Independently,
each drop is at most the negative variation on its gap. A 1-Lipschitz
function has total variation at most the interval length, so
\[
 T\le\operatorname{Var}^-(g;[p,1])
 \le\frac{1-p+g(p)-g(1)}2.
\]
This is the second bound. Their convex combination with weights
\(1-2\chi\) and \(2\chi\) proves \eqref{hardgap:blended-tail}.
\end{proof}

If there is just one positive gap, the barriers and
\eqref{hardgap:geometric-drop} give
\begin{align}
 \delta
 &\le\min\left\{p-\frac{1+q}{2},\ bq-ap\right\}
 \le\frac{(2b-a)p-b}{1+2b}
 \le C_1(a,b).\label{hardgap:one-gap}
\end{align}
The middle inequality follows by intersecting the decreasing and increasing
affine bounds in \(q\); the last follows from \(p\le1\).
With no positive gaps, the cost is zero.

\subsection{The two-collapse upper bound}
\label{hardgap:low-proof}

Assume \(b\le b_{\mathrm c}(a)\) and \(R(a,b)\le0\). Put
\[
 s=1+2b,\qquad D=(1-a)s^2,\qquad
 \chi=\frac{1+a-4ab+4b^2-2b}{D}.
\]
The numerator is \(1-2b+a+4b(b-a)>0\), while
\begin{equation}\label{hardgap:chi-sign}
 \chi-\frac12=\frac{R(a,b)}{2D}\le0.
\end{equation}
Thus \(0<\chi\le1/2\), and the blended tail bound is applicable.

Suppose there are at least two positive gaps, and denote the first two by
\((q_i,p_i)\), with \(u_i=g(q_i)\), \(v_i=g(p_i)\), and
\(\delta_i=u_i-v_i\), for \(i=1,2\). They obey
\begin{equation}\label{hardgap:low-constraints}
 \begin{aligned}
 u_1-bq_1&\le0,& ap_1-v_1&\le0,\\
 q_1-2p_1+2\delta_1&\le-1,&
 u_2-v_1-q_2+p_1&\le0,\\
 u_2-bq_2&\le0,& q_2-2p_2+2\delta_2&\le-1,\\
 p_2&\le1.&&
 \end{aligned}
\end{equation}
The fourth inequality uses Lipschitz continuity between \(p_1\) and
\(q_2\), regardless of any omitted zero-drop gaps.

In the displayed order, multiply the seven constraints by
\begin{equation}\label{hardgap:low-weights}
 \begin{aligned}
 w_1&=\frac1s,&
 w_2&=\frac{1-2b}{(1-a)s},&
 w_3&=\frac b{s},\\
 w_4&=\frac{2b-a}{(1-a)s},&
 w_5&=\frac{1+2a-4b-2ab}{D},&
 w_7&=\frac{1-\chi}{2},\\
 w_8&=\frac12-\chi.&&&&
 \end{aligned}
\end{equation}
Every weight is nonnegative. The sign of \(w_5\) is exactly
\(b\le b_{\mathrm c}(a)\); all other signs follow from the parameter
range and \eqref{hardgap:chi-sign}. The subscripts are merely convenient
names for these seven coefficients.

The identities
\[
 \begin{gathered}
 w_3=bw_1,\quad w_1+2w_3=1,\quad
 aw_2-2w_3+w_4=0,\\
 w_2+2w_3+w_4=1,\quad
 w_7=w_4+bw_5,\quad w_4+w_5=\chi,\quad
 2w_7=1-\chi
 \end{gathered}
\]
check every coefficient in the weighted sum. Its left side is
\(\delta_1+\delta_2-p_2/2+\chi v_2\), and its right side is
\(-w_3-w_7+w_8=-w_3-\chi/2\). Therefore
\begin{equation}\label{hardgap:low-certificate}
 \delta_1+\delta_2-\frac{p_2}{2}+\chi v_2
 \le-\frac b{1+2b}-\frac\chi2.
\end{equation}
Apply \eqref{hardgap:blended-tail} to all gaps after \(p_2\). Using
\eqref{hardgap:drop-sum} gives
\begin{align*}
 \Phi_0(g)
 &\le\delta_1+\delta_2+\frac{1-p_2}{2}+\chi(v_2-a)\\
 &\le\frac12-\frac b{1+2b}-\left(a+\frac12\right)\chi
 =C_2(a,b).
\end{align*}
The zero- and one-positive-gap cases follow from
\eqref{hardgap:one-gap}, since
\begin{equation}\label{hardgap:cost-difference}
 C_2(a,b)-C_1(a,b)
 =\frac{(b-a)(1+2a-4b-2ab)}{(1-a)(1+2b)^2}\ge0.
\end{equation}
This proves the low-side upper bound for piecewise affine profiles.

\subsection{The single-collapse upper bound}
\label{hardgap:high-proof}

Now assume \(b\ge b_{\mathrm c}(a)\). In addition to the constraints
already established, the first two positive gaps satisfy
\begin{equation}\label{hardgap:between-gaps}
 v_1\le u_2.
\end{equation}
Indeed, between \(p_1\) and \(q_2\) the hard-component increases
are nonnegative and all intervening gaps have zero drop. Telescoping gives
\(u_2-v_1\ge0\). Pointwise monotonicity on that whole interval is
neither asserted nor needed.

Use the following eight constraints, in order:
\begin{equation}\label{hardgap:high-constraints}
 \begin{aligned}
 u_1-bq_1&\le0,\\
 \delta_1-p_1+q_1/2&\le-1/2,\\
 ap_1-v_1&\le0,\\
 \delta_2-p_2+q_2/2&\le-1/2,\\
 ap_2-v_2&\le0,\\
 u_2-v_1-q_2+p_1&\le0,\\
 p_2&\le1,\\
 v_1-u_2&\le0.
 \end{aligned}
\end{equation}
With \(D_1=(1+2a)(1+2b)\), multiply them by
\begin{equation}\label{hardgap:high-weights}
 \begin{aligned}
 \omega_1&=\frac1{1+2b},&
 \omega_2&=\frac{2b}{1+2b},\\
 \omega_3&=\frac{2b}{D_1},&
 \omega_4&=\frac{4b(1+a)}{D_1},\\
 \omega_5&=\frac{1+2a-2b}{D_1},&
 \omega_6&=\frac{2b(1+a)}{D_1},\\
 \omega_7&=\frac{(6+8a)b-(1+2a)^2}{2D_1},&
 \omega_8&=\frac{(2a+4)b-(2a+1)}{D_1}.
 \end{aligned}
\end{equation}
All are nonnegative. The last has precisely the required threshold.
The numerator of \(\omega_7\) increases with \(b\), and at
\(b=b_{\mathrm c}(a)\) it equals
\[
 \frac{(1+a)(1-2a)(1+2a)}{a+2}>0.
\]
The remaining signs follow from \(0<a<b\le1/2\).

Expanding the weighted sum of \eqref{hardgap:high-constraints} gives
exactly \(\delta_1+\delta_2-p_2/2\) on the left. On the right it
gives \(-\omega_2/2-\omega_4/2+\omega_7\). Thus the first tail
bound in \eqref{hardgap:tails} proves
\begin{align*}
 \Phi_0(g)
 &\le\delta_1+\delta_2+\frac{1-p_2}{2}\\
 &\le\frac12-\frac{\omega_2}{2}-\frac{\omega_4}{2}+\omega_7
 =C_1(a,b).
\end{align*}
The cases with fewer than two positive gaps follow from
\eqref{hardgap:one-gap}. This proves the second upper bound for piecewise
affine profiles.

\subsection{Finite starts, uniform chain length, and stability}
\label{hardgap:compactness}

We now justify the finite-chain statement and extend the proof to all
continuous profiles. First, \(\Phi_\zeta(g)\) is nondecreasing as
\(\zeta\downarrow0\). To compare starts \(n_1<n_2<1\), truncate
a chain from \(n_2\) at its first crossing of \(n_1\), replacing the
crossing edge by its subinterval from \(n_1\). The subinterval is
admissible and has no greater cost. Conversely, maximal admissible edges
bridge \(n_2\) to \(n_1\), shortening the last edge if necessary,
with total cost at most \(n_2-n_1\). Hence
\begin{equation}\label{hardgap:start-error}
 0\le\Phi_0(g)-\Phi_\zeta(g)\le\zeta.
\end{equation}
In particular, the proved limiting upper estimates bound every finite-start
infimum.

Any finite admissible chain can be compressed without increasing cost.
For \(l<m<n\),
\begin{equation}\label{hardgap:merge}
 c_g(l,n)\le c_g(l,m)+c_g(m,n).
\end{equation}
If \(A=\min_{[l,m]}g\), \(B=\min_{[m,n]}g\), the difference
between the right and left sides is \(g(m)-\max(A,B)\ge0\).
Merge adjacent edges whenever their union remains admissible. When no
merge is possible, three consecutive endpoints satisfy
\[
 2n_i-n_{i+2}>1,\qquad
 1-n_{i+2}>2(1-n_i).
\]
The complementary depth more than doubles every two edges. A start
\(n_0=1-\zeta\) therefore yields at most
\(2\lceil\log_2(1/\zeta)\rceil+O(1)\) retained edges.

Take finite chains with costs approaching the finite-start infimum and
compress them. Pad their endpoint lists by repetitions to a common bounded
length and extract a convergent subsequence. Admissibility is closed, and
interval minima vary continuously with the endpoints of their intervals.
The limiting chain attains the infimum; removing repeated endpoints makes
it strictly decreasing. Thus the desired upper cost is achieved by an
actual chain of uniformly bounded length.

For a continuous 1-Lipschitz profile, first normalize its endpoint as above,
then interpolate on partitions of mesh tending to zero, containing
\(0,n_0,1\). The interpolants preserve both linear barriers, the
Lipschitz constant, and the value \(g(1)=a\). The preceding argument
provides chains of uniformly bounded length with the asserted cost for
each interpolant. Their endpoint lists have a convergent subsequence, and
the interpolants converge uniformly. Interval minima and costs pass to the
limit, as does admissibility. This gives the asserted chain for the
continuous profile. Equation~\eqref{hardgap:clip} returns it to the
original profile if normalization was needed.

For grid rounding, rescale back to \([0,N]\). Round every endpoint
down to a multiple of \(T\). If \(m\ge2n-N\) and \(N/T\)
is integral, then
\[
 \left\lfloor\frac mT\right\rfloor
 \ge\left\lfloor\frac{2n-N}{T}\right\rfloor
 \ge2\left\lfloor\frac nT\right\rfloor-\frac NT.
\]
Thus admissibility is preserved. Endpoints move by at most \(T\), and
the interval minimum of a 1-Lipschitz function changes by at most \(T\).
Each cost changes by at most \(2T\); discard repetitions.

If the barriers hold with additive error \(\epsilon N\), clamp the
profile between \(ax\) and \(bx\) first. This preserves the Lipschitz
constant and changes the profile uniformly by at most \(\epsilon N\).
Each edge cost changes by at most \(2\epsilon N\). Normalize the
clamped profile afterwards; that step only increases edge costs and does
not add an error. The resulting total error is
\(O(K\epsilon N+KT)\), proving all claims concerning finite chains
and perturbations in Theorem~\ref{hardgap:profile}.

\subsection{A matching two-collapse witness}
\label{hardgap:two-witness}

We give the full hard-set calculation for the lower bound. Suppose
\(0<a<b\le b_{\mathrm c}(a)\), and define
\[
 k=2b-a,\qquad q_0=\frac b{k},\qquad
 Q(v)=\frac{2(1+a)v-1}{1+2b}.
\]
Set
\[
 q=Q(1)=\frac{1+2a}{1+2b},\qquad
 t=\frac{(1-b)q}{1-a},\qquad r=Q(t),
\]
and consider
\begin{equation}\label{hardgap:two-profile}
 g(x)=\min\{bx,\ at+|x-t|,\ 1+a-x\}.
\end{equation}
This is 1-Lipschitz. Each function in the minimum majorizes \(ax\):
for the middle one, its difference from \(ax\) is
\((1+a)(t-x)\) when \(x\le t\), and \((1-a)(x-t)\) when
\(x\ge t\). The last function has difference \((1+a)(1-x)\).
Thus the profile has both barriers, and \(g(1)=a\).

The parameter restriction gives \(q_0\le t<q<1\); the first inequality
is equivalent to \(1+2a-4b-2ab\ge0\). Put
\[
 L=\frac{(1+a)t}{1+b},\qquad P=\frac{1+a}{1+b}.
\]
Then \(L<t<q<P<1\), and the consecutive affine pieces are
\begin{equation}\label{hardgap:two-pieces}
 g(x)=\begin{cases}
 bx,&0\le x\le L,\\
 (1+a)t-x,&L\le x\le t,\\
 x-(1-a)t,&t\le x\le q,\\
 bx,&q\le x\le P,\\
 1+a-x,&P\le x\le1.
 \end{cases}
\end{equation}
The increasing cone after \(t\) meets \(bx\) exactly at \(q\),
because \((1-a)t=(1-b)q\).

The hard set has exactly the nondegenerate components
\begin{equation}\label{hardgap:two-hard}
 [0,r]\quad\text{and}\quad[t,q],
\end{equation}
together with the isolated endpoint one. To check this, first note
\[
 r-L=\frac{(1+a)t-(1+b)}{(1+b)(1+2b)}<0,
 \qquad \frac{1+r}{2}\le t,
\]
where the second inequality is equivalent to \(t\ge q_0\).
Also
\[
 br=(1+a)t-\frac{1+r}{2},\qquad
 br-at=\frac{kt-b}{1+2b}\ge0.
\]
For \(x\le r\), the interval \([x,(1+x)/2]\) either stays on the
initial increasing segment or enters its following descent with endpoint
value at least \(bx\). Thus these points are hard. For \(r<x<L\),
once its future endpoint is on that descent it has value below \(bx\);
if the interval reaches \(t\), its value \(at<bx\) gives the same
conclusion. The relevant endpoint is already on or beyond the descent at
\(x=r\), since
\[
 \frac{1+r}{2}-L
 =\frac{b(1+b-(1+a)t)}{(1+b)(1+2b)}\ge0.
\]
For \(L\le x<t\), immediate strict decrease rules out hardness.
There are consequently no hard points in \((r,t)\).

For \(t\le x\le q\), the profile initially increases. If its future
interval enters the final descent, the difference between the final-descent
value at \((1+x)/2\) and \(g(x)\) equals
\[
 1+a-\frac{1+x}{2}-\bigl(x-(1-a)t\bigr)
 =\frac32(q-x)\ge0.
\]
Thus these points are hard. For \(q<x\le P\), the same future
endpoint lies on the final descent and has value strictly below \(bx\).
Indeed,
\[
 \frac{1+q}{2}-P
 =\frac{b(b-a)}{(1+2b)(1+b)}>0.
\]
For \(P\le x<1\), strict decrease again rules out hardness.
This proves \eqref{hardgap:two-hard}, including the boundary case
\(t=q_0\), where the first gap has zero drop.

The identity \eqref{hardgap:identity} now gives
\[
 \Phi_0(g)=br+(bq-at)-a=C_2(a,b).
\]
Thus the low-side upper estimate is attained by a concrete piecewise affine
profile wherever its hypotheses hold.

\subsection{A matching single-collapse witness}
\label{hardgap:one-witness}

For all \(0<a<b\le1\), the profile
\[
 g(x)=\min\{bx,1+a-x\}
\]
is admissible and endpoint-normalized. Its only nondegenerate hard component
is
\[
 \left[0,\frac{1+2a}{1+2b}\right].
\]
Indeed, on the initial increasing segment the boundary is the solution of
\(bx=1+a-(1+x)/2\); beyond it the future descent has smaller value,
and the final decreasing segment has no hard points except the endpoint
one. Therefore
\[
 \Phi_0(g)=b\frac{1+2a}{1+2b}-a=C_1(a,b).
\]
This proves sharpness on the high-side domain and completes the proof of
Theorem~\ref{hardgap:profile}. By \eqref{hardgap:cost-difference}, the two
sharp formulas agree at \(b=b_{\mathrm c}(a)\).
